%%%%%%%%%%%%%%%%%%%%%%%%%%%%%%%%%%%%%%%%%%%%%%%%%%%%%%%%%%%%%%%%%%%%%%%%%%%%%%%%
% Obxectivo: Lista de termos empregados no documento,                          %
%            xunto cos seus respectivos significados.                          %
%%%%%%%%%%%%%%%%%%%%%%%%%%%%%%%%%%%%%%%%%%%%%%%%%%%%%%%%%%%%%%%%%%%%%%%%%%%%%%%%

\newglossaryentry{backend}{
  name=backend,
  description={Componente servidor responsable de exponer la lógica de negocio, gestionar la persistencia y ofrecer una interfaz de comunicación a otros clientes o sistemas.}
}

\newglossaryentry{contenedor}{
  name=contenedor,
  description={Unidad de empaquetado y ejecución que incluye una aplicación y sus dependencias para facilitar despliegues reproducibles.}
}

\newglossaryentry{infraestructura-como-codigo}{
  name=infraestructura como código,
  description={Práctica consistente en definir recursos de infraestructura mediante ficheros versionables y automatizables.}
}

\newglossaryentry{stateless}{
  name=stateless,
  description={Propiedad de un servicio que no conserva estado de sesión en memoria entre peticiones y puede escalarse más fácilmente detrás de un balanceador.}
}

\newglossaryentry{runbook}{
  name=runbook,
  description={Procedimiento operativo documentado que describe cómo ejecutar, diagnosticar o recuperar una operación técnica concreta.}
}

\newglossaryentry{rollback}{
  name=rollback,
  description={Vuelta controlada a una versión previa conocida de una aplicación o configuración cuando un cambio nuevo falla.}
}
